\section{アプローチ05（オブジェクト検出とポーズ検出のハイブリッド手法）}

\subsection{手法の概要}
アプローチ04で判明した「高解像度画像全体に対する直接的なポーズ推定の計算コスト」および「密集環境でのキーポイント交錯」を解決するため、アプローチ05では一次検出として信頼度閾値を低めに設定したオブジェクト検出を行い、その切り出し領域（ROI）に対してポーズ推定を適用するハイブリッド検出パイプラインを導入した。

後ろ姿対策と体の向き検出を兼ねるため、ポーズ推定モデルには MoveNet Lightning を採用した。処理の効率化とエッジ環境（Webブラウザ）におけるレスポンス向上を図るため、以下の構成および並行・メモリ最適化を実装した。

\begin{itemize}
    \item \textbf{切り出しと所有権移転}：画像の切り出しは、固定サイズの非表示キャンバスおよび \texttt{createImageBitmap} を用いてバウンディングボックスの範囲を指定して実行し、Web Worker 間では転送（Transferable Objects）による所有権移転を活用して不要なデータコピーやメモリアロケーションを削減した。
    \item \textbf{サイズ同期}：Worker Manager から各 Worker へ初回起動時に元画像のサイズを伝達し、内部で保持する固定キャンバスの解像度を事前に決定させた。
    \item \textbf{切り出しマージン}：人物が密集している状況においても隣接オブジェクトの混入を避けるため、画像切り出し時の余白（マージン）は一切付与しない仕様とした。
    \item \textbf{固定Workerプール}：多数の候補に対する処理遅延やリソース暴走を防ぐため、並列実行を行う Worker プール数を固定（4コアを目安）とし、全体の要求を制御した。
    \item \textbf{非同期パイプライン}：源流（メイン処理）の \texttt{Promise.all} にてバウンディングボックス列挙で得られた候補を一括生成し、同一の生成済みキャンバスからビットマップで切り出したうえで Worker Manager へ同時投入した。さらに、得られた処理結果（Worker からの返却データ）はメイン側で元画像の座標系へと逆変換・復元した。
\end{itemize}

\subsection{結果およびアプローチ04との比較・課題}

\subsubsection{アプローチ04との比較}
画面全体へ直接ポーズ推定を適用していたアプローチ04と比較し、一次検出で絞り込んだ領域のみに MoveNet Lightning を適用することで計算コストが劇的に削減された。また、所有権移転や Worker プールの活用により、UIスレッドをブロックすることなく実用に耐えうる処理スピードと低信頼度フィルタリングの精度を両立させた。

\subsubsection{判明した課題}
\begin{enumerate}
    \item \textbf{小サイズバウンディングボックスにおける推定精度の低下}：
    一次検出で列挙された人物領域が小さすぎる場合、MoveNet Lightning の入力要求サイズを満たせず、キーポイントの検出精度が著しく低下する（または検出自体に失敗する）問題が判明した。
    \item \textbf{重複するバウンディングボックスの処理負荷}：
    低信頼度で手当たり次第に候補を抽出するため、同一人物に対して複数の領域が切り出され、Worker プールへ無駄なタスクが集中する問題が浮彫りとなった。
\end{enumerate}
